\documentclass[../../main.tex]{subfiles}
\begin{document}
\begin{oframed}
  \subfile{problem}
\end{oframed}

{\bf [解答]}
$e(\theta) = \cos\theta + i\sin\theta$ とする．$z = e(\alpha)$, $w = e(\beta)$ であり，
\begin{align}
  z+w
  &= e(\alpha) + e(\beta) \\
  &= (\cos\alpha + \cos\beta) + i(\sin\alpha + \sin\beta) \\
  &= 2 \cos\dfrac{\alpha+\beta}{2} \cos\dfrac{\alpha-\beta}{2} + 2 i \sin\dfrac{\alpha+\beta}{2} \cos\dfrac{\alpha-\beta}{2} \\
  &= 2 \cos\dfrac{\alpha-\beta}{2} e\left(\dfrac{\alpha+\beta}{2}\right) \label{eq:1}
\end{align}
と書ける．

$t = z+w$ として題意の条件式
\begin{align}
  |1 + z + w| \le 1 \iff |1 + t|^2 \le 1 \iff |t|^2 + t + \bar{t} \le 0 \label{eq:2}
\end{align}
に\cref{eq:1}を代入する．$|e(\theta)| = 1$ に注意して変形すると
\begin{align}
  & 4 \cos^2\dfrac{\alpha-\beta}{2}\left|e\left(\dfrac{\alpha+\beta}{2}\right)\right|^2 + 2 \cos\dfrac{\alpha-\beta}{2} \cdot 2 \cos\dfrac{\alpha+\beta}{2} \le 0 \\
  & 4 \cos^2\dfrac{\alpha-\beta}{2} + 2 \cos\dfrac{\alpha-\beta}{2} \cdot 2 \cos\dfrac{\alpha+\beta}{2} \le 0 \\
  & \cos\dfrac{\alpha-\beta}{2} \left(\cos\dfrac{\alpha-\beta}{2} + \cos\dfrac{\alpha+\beta}{2}\right) \le 0 \\
  \therefore
  & \cos\dfrac{\alpha-\beta}{2} \cos\dfrac{\alpha}{2} \cos\dfrac{\beta}{2} \le 0 \label{eq:3}
\end{align}
を得る．

以上から求める $(\alpha, \beta)$ は\cref{eq:3}を満たし，かつ $0 \le \alpha, \beta \le 2\pi$ なる $(\alpha, \beta)$ である．
$\cos\dfrac{\alpha}{2}, \cos\dfrac{\beta}{2}$の正負に注意して整理すると解は以下の通りである．
\begin{align}
  \begin{cases}
    \alpha = \pi \text{ の時} & 0 \le \beta \le 2\pi \\
    \beta = \pi \text{ の時} & 0 \le \alpha \le 2\pi \\
    0 \le \alpha < \pi \land 0 \le \beta < \pi \text{ の時}       & \cos\dfrac{\alpha-\beta}{2} \le 0 \\
    0 \le \alpha < \pi \land \pi < \beta \le 2\pi \text{ の時}    & \cos\dfrac{\alpha-\beta}{2} \ge 0 \\
    \pi < \alpha \le 2\pi \land 0 \le \beta < \pi \text{ の時}    & \cos\dfrac{\alpha-\beta}{2} \ge 0 \\
    \pi < \alpha \le 2\pi \land \pi < \beta \le 2\pi \text{ の時} & \cos\dfrac{\alpha-\beta}{2} \le 0
  \end{cases}
\end{align}

これを図示すると下図斜線部（境界含む）である．

\begin{figure}[htb]
  \centering
  \begin{tikzpicture}[scale=0.95, >=stealth, every node/.style={font=\small}]
    % [0, 2pi]^2 の基準枠
    \draw[densely dotted, thick, gray!70!black] (0, 0) rectangle ({2*pi}, {2*pi});

    % 2つの三角形領域（斜線）
    \fill[pattern=north east lines, pattern color=blue!65] (0, {pi}) -- ({pi}, {pi}) -- ({pi}, {2*pi}) -- cycle;
    \fill[pattern=north east lines, pattern color=blue!65] ({pi}, 0) -- ({pi}, {pi}) -- ({2*pi}, {pi}) -- cycle;

    % 座標軸
    \draw[->, thick, gray!80!black] (-0.4, 0) -- ({2*pi+0.7}, 0) node[right] {$\alpha$};
    \draw[->, thick, gray!80!black] (0, -0.4) -- (0, {2*pi+0.7}) node[above] {$\beta$};

    % 解の境界および直線 alpha=pi, beta=pi 上の線分（すべて含む）
    \draw[line width=1.4pt, blue!85!black] (0, {pi}) -- ({pi}, {2*pi});
    \draw[line width=1.4pt, blue!85!black] ({pi}, 0) -- ({2*pi}, {pi});
    \draw[line width=1.4pt, blue!85!black] ({pi}, 0) -- ({pi}, {2*pi});
    \draw[line width=1.4pt, blue!85!black] (0, {pi}) -- ({2*pi}, {pi});

    % 境界上の端点
    \node[circle, fill=blue!85!black, inner sep=1.3pt] at (0, {pi}) {};
    \node[circle, fill=blue!85!black, inner sep=1.3pt] at ({pi}, 0) {};
    \node[circle, fill=blue!85!black, inner sep=1.3pt] at ({pi}, {2*pi}) {};
    \node[circle, fill=blue!85!black, inner sep=1.3pt] at ({2*pi}, {pi}) {};

    % 目盛りラベル
    \node[below left=2pt] at (0, 0) {$0$};
    \node[below=2pt] at ({pi}, 0) {$\pi$};
    \node[below=2pt] at ({2*pi}, 0) {$2\pi$};
    \node[left=2pt] at (0, {pi}) {$\pi$};
    \node[left=2pt] at (0, {2*pi}) {$2\pi$};
  \end{tikzpicture}
  \caption{\cref{eq:2}を満たす$(\alpha,\beta)$の範囲（斜線部と$\alpha=\pi$，$\beta=\pi$の太線部，境界含む）}
  \label{fig:alpha_beta_region}
\end{figure}

\end{document}
